\chapter{Tests et assurance qualite}

\section{Qualite backend}

Le backend integre nativement plusieurs leviers qualite:
\begin{itemize}[leftmargin=*]
  \item tests unitaires via Surefire;
  \item tests integration via Failsafe;
  \item couverture via JaCoCo;
  \item analyse statique via SpotBugs et Checkstyle (profil \texttt{analysis}).
\end{itemize}

Le pipeline GitLab automatise ces etapes avant la creation/push d'image Docker.

\section{Qualite frontend}

Le frontend Flutter peut etre verifie avec:
\begin{itemize}[leftmargin=*]
  \item \texttt{flutter analyze} pour l'analyse statique;
  \item \texttt{flutter test} pour tests unitaires/widgets;
  \item tests manuels sur les parcours critiques (auth, navigation role-based, chat, notifications).
\end{itemize}

\section{Qualite chatbot}

Pour le chatbot, trois niveaux de verification sont recommandes:
\begin{itemize}[leftmargin=*]
  \item tests fonctionnels de l'API \texttt{/api/chat};
  \item evaluation precision sur un jeu de questions de reference;
  \item suivi des scores similarity pour ajuster \texttt{MIN\_SIM} et \texttt{FUZZY\_MIN}.
\end{itemize}

\section{Scenario de validation bout en bout}

\begin{enumerate}[leftmargin=*]
  \item Connexion utilisateur Flutter et recuperation JWT.
  \item Chargement des donnees metier depuis le backend.
  \item Ouverture du module chat et verification des updates temps reel.
  \item Appel chatbot et verification du format de reponse.
  \item Verification des restrictions de role sur ecrans admin.
\end{enumerate}

\section{Indicateurs de suivi proposes}

\begin{table}[H]
\centering
\begin{tabular}{p{0.38\textwidth} p{0.48\textwidth}}
\toprule
\textbf{Indicateur} & \textbf{Objectif} \\
\midrule
Taux de succes API backend & Mesurer la stabilite des endpoints. \\
Temps de reponse chat temps reel & Valider la fluidite conversationnelle. \\
Precision chatbot (Top-1) & Evaluer la pertinence des reponses. \\
Taux erreurs UI critiques & Reduire les regressions fonctionnelles. \\
\bottomrule
\end{tabular}
\caption{KPIs qualite recommandes}
\end{table}
